%%%%%%%%%%%%%%%%
%%% num 9 %%%
%%%%%%%%%%%%%%%%

\Chapter{9}

\Verse{1}
{
	\hP{וַיְדַבֵּר יְהֹוָה אֶל מֹשֶׁה בְמִדְבַּר סִינַי בַּשָּׁנָה הַשֵּׁנִית לְצֵאתָם מֵאֶרֶץ מִצְרַיִם בַּחֹדֶשׁ הָרִאשׁוֹן לֵאמֹר׃}
}
{
	\eP{
		9 And YHWH had spoken to Moses in the wilderness of Sinai\\
		in the second year of their exodus from the land of Egypt,\\
		in the first month, saying, [image]\\
	}
}
{
%	\footnote{
%		had spoken. This is a rare occurrence in the Torah: telling an event out
%		of order. The book of Numbers begins with God speaking to Moses in the
%		second month of the year. Now it tells of God speaking to Moses back in
%		the first month of the year. Why? Because it tells of a case in which
%		some people had been unable to observe Passover, which falls in the
%		first month, and so God tells Moses in this flashback that the law in
%		that situation is that they should observe Passover one month later, on
%		the evening of the fourteenth day of the second month. Their observance
%		on that date thus belongs here in the sequence of the narrative. The
%		flashback is placed here to explain it. The point is that the Torah
%		adheres to chronological order in its story, and exceptional breaks in
%		that pattern such as this one occur very rarely and for specific
%		narrative purposes.
%	}
}

\Verse{2}
{
	\hP{וְיַעֲשׂוּ בְנֵי יִשְׂרָאֵל אֶת הַפָּסַח בְּמוֹעֲדוֹ׃}
}
{
	\eP{
		“And the children of Israel shall do the Passover at its\\
		appointed time. [image]\\
	}
}
{
}

\Verse{3}
{
	\hP{בְּאַרְבָּעָה עָשָׂר יוֹם בַּחֹדֶשׁ הַזֶּה בֵּין הָעַרְבַּיִם תַּעֲשׂוּ אֹתוֹ בְּמֹעֲדוֹ כְּכל חֻקֹּתָיו וּכְכל מִשְׁפָּטָיו תַּעֲשׂוּ אֹתוֹ׃}
}
{
	\eP{
		On the fourteenth day in this month ‘between the two\\
		evenings’ you shall do it at its appointed time. You shall\\
		do it according to all of its law and according to all of\\
		its judgments.” [image]\\
	}
}
{
%	\footnote{
%		between the two evenings. Understood to mean the early evening hours:
%		“dusk” or (The Hebrew has the dual ending, which produces the meaning of
%		being between two evenings.)
%	}
}

\Verse{4}
{
	\hP{וַיְדַבֵּר מֹשֶׁה אֶל בְּנֵי יִשְׂרָאֵל לַעֲשֹׂת הַפָּסַח׃}
}
{
	\eP{
		And Moses had spoken to the children of Israel to do the\\
		Passover, [image]\\
	}
}
{
}

\Verse{5}
{
	\hP{וַיַּעֲשׂוּ אֶת הַפֶּסַח בָּרִאשׁוֹן בְּאַרְבָּעָה עָשָׂר יוֹם לַחֹדֶשׁ בֵּין הָעַרְבַּיִם בְּמִדְבַּר סִינָי כְּכֹל אֲשֶׁר צִוָּה יְהֹוָה אֶת מֹשֶׁה כֵּן עָשׂוּ בְּנֵי יִשְׂרָאֵל׃}
}
{
	\eP{
		and they had done the Passover in the first month on the\\
		fourteenth day of the month “between the two evenings” in\\
		the wilderness of Sinai. According to all that YHWH had\\
		commanded Moses, so the children of Israel had done.\\
		[image]\\
	}
}
{
}

\Verse{6}
{
	\hP{וַיְהִי אֲנָשִׁים אֲשֶׁר הָיוּ טְמֵאִים לְנֶפֶשׁ אָדָם וְלֹא יָכְלוּ לַעֲשֹׂת הַפֶּסַח בַּיּוֹם הַהוּא וַיִּקְרְבוּ לִפְנֵי מֹשֶׁה וְלִפְנֵי אַהֲרֹן בַּיּוֹם הַהוּא׃}
}
{
	\eP{
		And there had been people who were impure by a human being\\
		and so had not been able to do the Passover on that day,\\
		and they had come forward in front of Moses and in front\\
		of Aaron on that day, [image]\\
	}
}
{
%	\footnote{
%		impure by a human being. See comment on Num 5:2.
%	}
}

\Verse{7}
{
	\hP{וַיֹּאמְרוּ הָאֲנָשִׁים הָהֵמָּה אֵלָיו אֲנַחְנוּ טְמֵאִים לְנֶפֶשׁ אָדָם לָמָּה נִגָּרַע לְבִלְתִּי הַקְרִיב אֶת קרְבַּן יְהֹוָה בְּמֹעֲדוֹ בְּתוֹךְ בְּנֵי יִשְׂרָאֵל׃}
}
{
	\eP{
		and those people had said to him, “We’re impure by a human\\
		being. Why should we be excluded so as not to make YHWH’s\\
		offering at its appointed time among the children of\\
		Israel?” [image]\\
	}
}
{
}

\Verse{8}
{
	\hP{וַיֹּאמֶר אֲלֵהֶם מֹשֶׁה עִמְדוּ וְאֶשְׁמְעָה מַה יְצַוֶּה יְהֹוָה לָכֶם׃}
}
{
	\eP{
		And Moses had said, “Stay, and let me hear what YHWH will\\
		command about you.” [image]\\
	}
}
{
}

\Verse{9}
{
	\hP{וַיְדַבֵּר יְהֹוָה אֶל מֹשֶׁה לֵּאמֹר׃}
}
{
	\eP{And YHWH had spoken to Moses, saying, [image]}
}
{
}

\Verse{10}
{
	\hP{דַּבֵּר אֶל בְּנֵי יִשְׂרָאֵל לֵאמֹר אִישׁ אִישׁ כִּי יִהְיֶה טָמֵא לָנֶפֶשׁ אוֹ בְדֶרֶךְ רְחֹקָה לָכֶם אוֹ לְדֹרֹתֵיכֶם וְעָשָׂה פֶסַח לַיהֹוָה׃}
}
{
	\eP{
		“Speak to the children of Israel, saying: Any man who will\\
		be impure by a person or on a faraway journey—among you or\\
		through your generations—and will do the Passover for\\
		YHWH, [image]\\
	}
}
{
}

\Verse{11}
{
	\hP{בַּחֹדֶשׁ הַשֵּׁנִי בְּאַרְבָּעָה עָשָׂר יוֹם בֵּין הָעַרְבַּיִם יַעֲשׂוּ אֹתוֹ עַל מַצּוֹת וּמְרֹרִים יֹאכְלֻהוּ׃}
}
{
	\eP{
		they shall do it in the second month on the fourteenth day\\
		‘between the two evenings.’ They shall eat it with\\
		unleavened bread and bitter herbs. [image]\\
	}
}
{
}

\Verse{12}
{
	\hP{לֹא יַשְׁאִירוּ מִמֶּנּוּ עַד בֹּקֶר וְעֶצֶם לֹא יִשְׁבְּרוּ בוֹ כְּכל חֻקַּת הַפֶּסַח יַעֲשׂוּ אֹתוֹ׃}
}
{
	\eP{
		They shall not leave any of it until morning, and they\\
		shall not break a bone of it. They shall do it according\\
		to all of the law of the Passover. [image]\\
	}
}
{
}

\Verse{13}
{
	\hP{וְהָאִישׁ אֲשֶׁר הוּא טָהוֹר וּבְדֶרֶךְ לֹא הָיָה וְחָדַל לַעֲשׂוֹת הַפֶּסַח וְנִכְרְתָה הַנֶּפֶשׁ הַהִוא מֵעַמֶּיהָ כִּי קרְבַּן יְהֹוָה לֹא הִקְרִיב בְּמֹעֲדוֹ חֶטְאוֹ יִשָּׂא הָאִישׁ הַהוּא׃}
}
{
	\eP{
		But a man who is pure and has not been on a journey and\\
		has failed to do the Passover: that person will be cut off\\
		from his people, because he did not make YHWH’s offering\\
		at its appointed time. That man shall bear his sin.\\
		[image]\\
	}
}
{
%	\footnote{
%		cut off from his people. On the meaning of this phrase, see the comment
%		on Gen 17:14.
%	}
}

\Verse{14}
{
	\hP{וְכִי יָגוּר אִתְּכֶם גֵּר וְעָשָׂה פֶסַח לַיהֹוָה כְּחֻקַּת הַפֶּסַח וּכְמִשְׁפָּטוֹ כֵּן יַעֲשֶׂה חֻקָּה אַחַת יִהְיֶה לָכֶם וְלַגֵּר וּלְאֶזְרַח הָאָרֶץ׃}
}
{
	\eP{
		And if an alien will reside with you and will do the\\
		Passover for YHWH: according to the law of the Passover\\
		and according to its required manner, so he shall do it.\\
		You shall have one law, for the alien and for the citizen\\
		of the land.” [image]\\
	}
}
{
}

\Verse{15}
{
	\hP{וּבְיוֹם הָקִים אֶת הַמִּשְׁכָּן כִּסָּה הֶעָנָן אֶת הַמִּשְׁכָּן לְאֹהֶל הָעֵדֻת וּבָעֶרֶב יִהְיֶה עַל הַמִּשְׁכָּן כְּמַרְאֵה אֵשׁ עַד בֹּקֶר׃}
}
{
	\eP{
		And on the day that the Tabernacle was set up, the cloud\\
		covered the Tabernacle of the tent of the Testimony, and\\
		in the evening it would be over the Tabernacle like the\\
		appearance of fire until morning. [image]\\
	}
}
{
%	\footnote{
%		the tent of the Testimony. This is the first occurrence of this phrase.
%		The usual phrase until now has been “Tent of Meeting.” The phrase here
%		may not refer to the entire Tent of Meeting, but rather to the inner
%		pavilion, the prket, which is associated with the Testimony (the tablets
%		of the Ten Commandments) four times (Exod 27:21; 30:6; Lev 24:3; Num
%		4:5) and is called explicitly “the pavilion of the Testimony” (Lev
%		24:3).
%	}
}

\Verse{16}
{
	\hP{כֵּן יִהְיֶה תָמִיד הֶעָנָן יְכַסֶּנּוּ וּמַרְאֵה אֵשׁ לָיְלָה׃}
}
{
	\eP{
		So it would be always: the cloud would cover it and the\\
		appearance of fire at night. [image]\\
	}
}
{
}

\Verse{17}
{
	\hP{וּלְפִי הֵעָלוֹת הֶעָנָן מֵעַל הָאֹהֶל וְאַחֲרֵי כֵן יִסְעוּ בְּנֵי יִשְׂרָאֵל וּבִמְקוֹם אֲשֶׁר יִשְׁכּן שָׁם הֶעָנָן שָׁם יַחֲנוּ בְּנֵי יִשְׂרָאֵל׃}
}
{
	\eP{
		And according to when the cloud was lifted from over the\\
		tent, after that the children of Israel would travel. And\\
		in a place where the cloud would tent, the children of\\
		Israel would camp there. [image]\\
	}
}
{
%	\footnote{
%		the cloud would tent. Forms of this Hebrew verb (yiškn) are frequently
%		taken to mean “to dwell”; but, when it refers to God or to the cloud,
%		this verb is a denominative from the noun miškn: Tabernacle. The deity
%		cannot be pictured as residing on the earth, but God rather is
%		understood to tent among humans, that is, to be present in association
%		with an impermanent, movable structure. For references, see the comment
%		on Exod 25:8.
%	}
}

\Verse{18}
{
	\hP{עַל פִּי יְהֹוָה יִסְעוּ בְּנֵי יִשְׂרָאֵל וְעַל פִּי יְהֹוָה יַחֲנוּ כּל יְמֵי אֲשֶׁר יִשְׁכֹּן הֶעָנָן עַל הַמִּשְׁכָּן יַחֲנוּ׃}
}
{
	\eP{
		By YHWH’s word the children of Israel would travel, and by\\
		YHWH’s word they would camp. All the days that the cloud\\
		would tent over the Tabernacle they would camp. [image]\\
	}
}
{
}

\Verse{19}
{
	\hP{וּבְהַאֲרִיךְ הֶעָנָן עַל הַמִּשְׁכָּן יָמִים רַבִּים וְשָׁמְרוּ בְנֵי יִשְׂרָאֵל אֶת מִשְׁמֶרֶת יְהֹוָה וְלֹא יִסָּעוּ׃}
}
{
	\eP{
		And when the cloud would extend many days over the\\
		Tabernacle, then the children of Israel kept YHWH’s charge\\
		and would not travel. [image]\\
	}
}
{
}

\Verse{20}
{
	\hP{וְיֵשׁ אֲשֶׁר יִהְיֶה הֶעָנָן יָמִים מִסְפָּר עַל הַמִּשְׁכָּן עַל פִּי יְהֹוָה יַחֲנוּ וְעַל פִּי יְהֹוָה יִסָּעוּ׃}
}
{
	\eP{
		And there were times that the cloud would be over the\\
		Tabernacle for a number of days: by YHWH’s word they would\\
		camp, and by YHWH’s word they would travel. [image]\\
	}
}
{
}

\Verse{21}
{
	\hP{וְיֵשׁ אֲשֶׁר יִהְיֶה הֶעָנָן מֵעֶרֶב עַד בֹּקֶר וְנַעֲלָה הֶעָנָן בַּבֹּקֶר וְנָסָעוּ אוֹ יוֹמָם וָלַיְלָה וְנַעֲלָה הֶעָנָן וְנָסָעוּ׃}
}
{
	\eP{
		And there were times that the cloud would be there from\\
		evening until morning, and the cloud was lifted in the\\
		morning, and then they would travel; or for a day and a\\
		night, and the cloud was lifted, and then they traveled.\\
		[image]\\
	}
}
{
}

\Verse{22}
{
	\hP{אוֹ יֹמַיִם אוֹ חֹדֶשׁ אוֹ יָמִים בְּהַאֲרִיךְ הֶעָנָן עַל הַמִּשְׁכָּן לִשְׁכֹּן עָלָיו יַחֲנוּ בְנֵי יִשְׂרָאֵל וְלֹא יִסָּעוּ וּבְהֵעָלֹתוֹ יִסָּעוּ׃}
}
{
	\eP{
		Whether two days or a month or a year, when the cloud\\
		would extend long over the Tabernacle, to tent over it,\\
		the children of Israel would camp and would not travel,\\
		and when it was lifted they would travel. [image]\\
	}
}
{
}

\Verse{23}
{
	\hP{עַל פִּי יְהֹוָה יַחֲנוּ וְעַל פִּי יְהֹוָה יִסָּעוּ אֶת מִשְׁמֶרֶת יְהֹוָה שָׁמָרוּ עַל פִּי יְהֹוָה בְּיַד מֹשֶׁה׃}
}
{
	\eP{
		By YHWH’s word they would camp, and by YHWH’s word they\\
		would travel. They kept YHWH’s charge, by YHWH’s word\\
		through Moses’ hand. [image]\\
	}
}
{
}
